
Calculons maintenant le second membre de la formule de Riemann--Roch. Pour tout @@M0@@, on a :
@@M1@@
d'après (V (6.2.1)). Par suite :
@@M2@@
Le second membre de la formule de Riemann--Roch est donc :
@@M3@@
Comme @@M4@@ est la classe dans @@M5@@ de @@M6@@, @@M7@@ est, pour @@M8@@, la classe dans
@@M9@@ de @@M10@@. On a donc :
@@M11@@
Rappelons enfin que les éléments @@M12@@ sont linéairement dépendants sur @@M13@@ (VI 5.6), et vérifient l'équation :
@@M14@@
et que @@M15@@ sont linéairement indépendants sur @@M16@@.

On peut alors écrire la série @@M17@@ de façon unique comme combinaison linéaire à coefficients dans
@@M18@@ des éléments @@M19@@. L'image par @@M20@@ de @@M21@@ est alors, d'après (5.4.4) et 3.7, son coefficient sur @@M22@@. Compte tenu de (5.4.1), le théorème de Riemann--Roch pour @@M23@@ est donc équivalent à :
@@M24@@
Ce coefficient est le résidu de la forme différentielle
@@M25@@
